% GENERATED FILE. Edit canonical lessons, then run npm run generate:books.
% canonical-source-hash: fnv1a64:1fcfb3433bdba296
% canonical-lessons: BN-C271-nouka, BN-W271-letter-au-kar, BN-C271-tikit-deben, BN-C271-hariye-gechhe, BN-R271-at-the-station

\chapter{Tickets, and When Things Go Wrong}
\label{ch:bn-tickets-and-when-things-go-wrong}
\hlchaptermodality{271}

\begin{chapteropening}
\textbf{By the end of this chapter:} \emph{I can buy a ticket (ekṭā ṭikiṭ deben?), ask the fare and the time, report a lost bag or a missed train (hāriye gechhe, chheṛe gechhe), and read and write \bn{ৌ} in \bn{নৌকা}.}

Complete BN-R271-at-the-station: buy a ticket and report a problem.
\end{chapteropening}

\section[\texorpdfstring{\bn{নৌকা} (noukā)}{নৌকা (noukā)}]{\bn{নৌকা} (noukā) — a boat}

\label{lesson:BN-C271-nouka}



\begin{quote}
Before the new one: ask \textquotedblleft{}How do I get to the station?\textquotedblright{} (\emph{sṭeshône kībhābe jābô?})
\end{quote}

\subsection*{You'll want to know: \bn{নৌকা}}
\textbf{\bn{নৌকা}} — \emph{noukā} — \textquotedblleft{}a boat\textquotedblright{}.

\begin{quote}
\textbf{\bn{আমরা} \bn{নৌকায়} \bn{নদী} \bn{পার} \bn{হলাম।}}
\end{quote}

\begin{quote}
\emph{āmrā noukāy nôdī pār hôlām.} — \textquotedblleft{}We crossed the river by boat.\textquotedblright{}
\end{quote}

\textbf{\bn{নৌকায়}}, with \textbf{-\bn{য়}}, is \textquotedblleft{}in a boat, by boat\textquotedblright{}, as \textbf{\bn{বাসে}} is \textquotedblleft{}by bus\textquotedblright{}. A \textbf{\bn{নৌকা}} is small, rowed or sailed on a river; the big ship is \textbf{\bn{জাহাজ}}, which you know.

\begin{cousinweb}
\textbf{\bn{নৌকা}} (\emph{noukā}) is a tatsama: Sanskrit \emph{naukā}, which Monier-Williams gives as \textquotedblleft{}a small boat or ship\textquotedblright{}. Beside it Monier-Williams lists \emph{nau}, \textquotedblleft{}a ship, boat\textquotedblright{}.
\end{cousinweb}

\subsection*{Your turn}
Say these aloud:

\begin{itemize}
\raggedright
  \item \emph{noukā}
  \item \textquotedblleft{}We crossed the river by boat\textquotedblright{} — \emph{āmrā noukāy nôdī pār hôlām}
  \item \textquotedblleft{}A boat is small; a ship is big\textquotedblright{} — \emph{noukā chhoṭo; jāhāj bôṛo}
\end{itemize}

\begin{tcolorbox}[breakable,colback=teal!4,colframe=teal!35!black,title={Before you move on}]
What does \textbf{\bn{নৌকা}} mean? (\textquotedblleft{}A boat\textquotedblright{}.) \textquotedblleft{}By boat\textquotedblright{}? (\textbf{\bn{নৌকায়}}.) The big ship? (\textbf{\bn{জাহাজ}}.)
\end{tcolorbox}
\section[\texorpdfstring{\bn{ৌ} (-ou)}{ৌ (-ou)}]{\bn{ৌ} (-ou) — the sign around the \bn{ন} of \bn{নৌকা}}

\label{lesson:BN-W271-letter-au-kar}



\begin{quote}
Before the new one: say \textbf{\bn{নৌকা}}, \textquotedblleft{}a boat\textquotedblright{}. Then write \textbf{\bn{নো}} (\emph{no}): \textbf{\bn{ে}} in front of the \textbf{\bn{ন}}, \textbf{\bn{া}} after it.
\end{quote}

\subsection*{Script you'll notice: \bn{ৌ}}
\textbf{\bn{ৌ}} — \emph{-ou}, the vowel of \textbf{\bn{নৌকা}}. Like \textbf{\bn{ো}}, it is written on both sides of its consonant: one part in front, one after.

Set the print of \textbf{\bn{নৌ}} beside \textbf{\bn{নো}}. In front of the \textbf{\bn{ন}}, both have \textbf{\bn{ে}}. After it, \textbf{\bn{নো}} has an upright stroke like \textbf{\bn{া}}. \textbf{\bn{নৌ}} has an upright stroke too, but its top stands above the headline in a hook, bent over to the left, its tip at the upper left.

\subsection*{Writing: \bn{ৌ} — copy the printed shape}
Put your pen beside \textbf{\bn{নৌ}} and copy the shape you can see, slowly and large: the \textbf{\bn{ে}} in front, the \textbf{\bn{ন}}, then the stroke after it with its hook above the headline. Check the copy against \textbf{\bn{নো}}: only the hook differs.

\begin{quote}
This book does not yet tell you where to start \bn{ৌ} or which way the pen goes. That is not written down here until it can be written down with a source. Copying the printed shape needs no such source, and it is how the shape becomes yours in the meantime.
\end{quote}

\subsection*{Your turn}
\begin{itemize}
\raggedright
  \item \emph{Find:} \bn{ৌ} around the \bn{ন} of \textbf{\bn{নৌকা}}
  \item \emph{Write it:} \textbf{\bn{নো}}, then \textbf{\bn{নৌ}}, from the print
  \item \emph{Write it:} \textbf{\bn{নৌকা}}, then \textbf{\bn{নৌকায়}}
  \item \emph{Write it:} \textbf{\bn{দূর}}, from memory
\end{itemize}

\begin{tcolorbox}[breakable,colback=teal!4,colframe=teal!35!black,title={Before you move on}]
Which sign is around the \textbf{\bn{ন}} of \textbf{\bn{নৌকা}}? (\textbf{\bn{ৌ}}.) How is it unlike \textbf{\bn{ো}}? (The stroke after the consonant has a hook above the headline.)
\end{tcolorbox}
\section[\texorpdfstring{\bn{একটা} \bn{টিকিট} \bn{দেবেন}? (ekṭā ṭikiṭ deben?)}{একটা টিকিট দেবেন? (ekṭā ṭikiṭ deben?)}]{\bn{একটা} \bn{টিকিট} \bn{দেবেন}? (ekṭā ṭikiṭ deben?) — at the ticket window}

\label{lesson:BN-C271-tikit-deben}



\begin{quote}
Say \textquotedblleft{}We crossed the river by boat\textquotedblright{}. (\emph{āmrā noukāy nôdī pār hôlām}.) Then \textbf{\bn{টিকিট}}, \textquotedblleft{}a ticket\textquotedblright{}.
\end{quote}

\subsection*{You'll want to know: asking for a ticket}
\begin{quote}
\textbf{\bn{একটা} \bn{টিকিট} \bn{দেবেন}?}
\end{quote}

\begin{quote}
\emph{ekṭā ṭikiṭ deben?} — \textquotedblleft{}Could I have a ticket, please?\textquotedblright{}
\end{quote}

Bengali often asks politely with the respectful future, \textquotedblleft{}will you give?\textquotedblright{}, where English says \textquotedblleft{}could I have\textquotedblright{}. \textbf{\bn{দেবেন}} is \textbf{\bn{দেওয়া}} in that form. To say where to, put the place in the \textbf{-\bn{র}} form: \textbf{\bn{গ্রামের} \bn{একটা} \bn{টিকিট}}, \textquotedblleft{}a ticket to the village\textquotedblright{}. To ask for a seat: \textbf{\bn{জানালার} \bn{পাশে} \bn{একটা} \bn{আসন} \bn{পাব}?} (\emph{jānālār pāshe ekṭā āsôn pābô?}), \textquotedblleft{}can I get a seat by the window?\textquotedblright{}.

\subsection*{You'll want to know: the fare and the time}
\begin{quote}
\textbf{\bn{ভাড়া} \bn{কত}? \bn{ট্রেন} \bn{কখন} \bn{ছাড়বে}?}
\end{quote}

\begin{quote}
\emph{bhāṛā kôtô? ṭren kôkhôn chhāṛbe?} — \textquotedblleft{}What's the fare? When does the train leave?\textquotedblright{}
\end{quote}

\textbf{\bn{কত}}, \textquotedblleft{}how much\textquotedblright{}, asks the fare. \textbf{\bn{কখন}}, \textquotedblleft{}when\textquotedblright{}, the third of chapter thirty's \textbf{\bn{যখন}}, \textbf{\bn{তখন}}, \textbf{\bn{কখন}}, asks the time. \textbf{\bn{ছাড়বে}}, \textquotedblleft{}will leave\textquotedblright{}, is \textbf{\bn{ছাড়া}}, \textquotedblleft{}to leave, to let go\textquotedblright{}, in the future.

\subsection*{Your turn}
Say these aloud:

\begin{itemize}
\raggedright
  \item \textquotedblleft{}Could I have a ticket to the village, please?\textquotedblright{} — \emph{grāmer ekṭā ṭikiṭ deben?}
  \item \textquotedblleft{}What's the fare?\textquotedblright{} — \emph{bhāṛā kôtô?}
  \item \textquotedblleft{}When does the train leave?\textquotedblright{} — \emph{ṭren kôkhôn chhāṛbe?}
\end{itemize}

\begin{tcolorbox}[breakable,colback=teal!4,colframe=teal!35!black,title={Before you move on}]
The polite \textquotedblleft{}could I have\textquotedblright{}? (\textbf{\bn{দেবেন}?}) \textquotedblleft{}How much is the fare?\textquotedblright{} (\textbf{\bn{ভাড়া} \bn{কত}?}) \textquotedblleft{}When will it leave?\textquotedblright{} (\textbf{\bn{কখন} \bn{ছাড়বে}?})
\end{tcolorbox}
\section[\texorpdfstring{\bn{আমার} \bn{ব্যাগ} \bn{হারিয়ে} … (āmār byāg hāriye …)}{আমার ব্যাগ হারিয়ে … (āmār byāg hāriye …)}]{\bn{আমার} \bn{ব্যাগ} \bn{হারিয়ে} \bn{গেছে} (āmār byāg hāriye gechhe) — when something goes wrong}

\label{lesson:BN-C271-hariye-gechhe}



\begin{quote}
Ask \textquotedblleft{}When does the train leave?\textquotedblright{} (\emph{ṭren kôkhôn chhāṛbe?}) Then say \textbf{\bn{হারিয়ে} \bn{যাওয়া}}, \textquotedblleft{}to get lost\textquotedblright{}, which you know.
\end{quote}

\subsection*{You'll want to know: something is lost}
\begin{quote}
\textbf{\bn{আমার} \bn{মানিব্যাগ} \bn{হারিয়ে} \bn{গেছে।}}
\end{quote}

\begin{quote}
\emph{āmār mānibyāg hāriye gechhe.} — \textquotedblleft{}My wallet has gone missing.\textquotedblright{}
\end{quote}

\textbf{\bn{গেছে}} is a form of \textbf{\bn{যাওয়া}} that says something happened and still matters now: the wallet went, and it is still gone. This book teaches it only in these set phrases. Whose it is goes in the \textbf{-\bn{র}} form: \textbf{\bn{আমার} \bn{ব্যাগ}}, \textbf{\bn{মীরার} \bn{পাসপোর্ট}}. Then what you will do: \textbf{\bn{পুলিশকে} \bn{জানাব}} (\emph{pulishke jānābô}), \textquotedblleft{}I'll tell the police\textquotedblright{}.

\subsection*{You'll want to know: the train has left}
\textbf{\bn{ছাড়া}}, \textquotedblleft{}to leave\textquotedblright{}, which you know, makes the same phrase: \textbf{\bn{ছেড়ে} \bn{গেছে}}, \textquotedblleft{}has left\textquotedblright{}.

\begin{quote}
\textbf{\bn{আমার} \bn{ট্রেন} \bn{ছেড়ে} \bn{গেছে।} \bn{পরের} \bn{ট্রেন} \bn{কখন}?}
\end{quote}

\begin{quote}
\emph{āmār ṭren chheṛe gechhe. pôrer ṭren kôkhôn?} — \textquotedblleft{}My train has left. When is the next train?\textquotedblright{}
\end{quote}

\subsection*{Your turn}
Say these aloud:

\begin{itemize}
\raggedright
  \item \textquotedblleft{}My passport has gone missing\textquotedblright{} — \emph{āmār pāsporṭ hāriye gechhe}
  \item \textquotedblleft{}I'll tell the police\textquotedblright{} — \emph{pulishke jānābô}
  \item \textquotedblleft{}My train has left. When is the next train?\textquotedblright{} — \emph{āmār ṭren chheṛe gechhe. pôrer ṭren kôkhôn?}
\end{itemize}

\begin{tcolorbox}[breakable,colback=teal!4,colframe=teal!35!black,title={Before you move on}]
\textquotedblleft{}My bag has gone missing\textquotedblright{}? (\textbf{\bn{আমার} \bn{ব্যাগ} \bn{হারিয়ে} \bn{গেছে}}.) \textquotedblleft{}The train has left\textquotedblright{}? (\textbf{\bn{ট্রেন} \bn{ছেড়ে} \bn{গেছে}}.) What does \textbf{\bn{গেছে}} add? (It happened, and it still matters now.)
\end{tcolorbox}
\section[(review)]{Review — at the station}

\label{lesson:BN-R271-at-the-station}



\begin{quote}
Ask for a ticket. (\emph{ekṭā ṭikiṭ deben?}) Then say \textquotedblleft{}My wallet has gone missing\textquotedblright{}. (\emph{āmār mānibyāg hāriye gechhe}.)
\end{quote}

\subsection*{Your turn}
\begin{itemize}
\raggedright
  \item \emph{Say it:} Could I have a ticket to the village, please? What's the fare?
  \item \emph{Say it:} When will the boat leave?
  \item \emph{Say it:} Excuse me, my bag has gone missing. I was sitting on the platform.
  \item \emph{Say it:} Our train has left. That's why we have to take the bus.
  \item \emph{Say it:} How do I get to the airport? How far is it?
  \item \emph{Write it:} \textbf{\bn{নৌকা}}, then \textbf{\bn{দূর}}
\end{itemize}

\begin{tcolorbox}[breakable,colback=teal!4,colframe=teal!35!black,title={Before you move on}]
The polite ask? (\textbf{\bn{দেবেন}?}) \textquotedblleft{}Has gone missing\textquotedblright{}? (\textbf{\bn{হারিয়ে} \bn{গেছে}}.) \textquotedblleft{}Has left\textquotedblright{}? (\textbf{\bn{ছেড়ে} \bn{গেছে}}.) Which sign is around the \textbf{\bn{ন}} of \textbf{\bn{নৌকা}}? (\textbf{\bn{ৌ}}.)
\end{tcolorbox}
